% Muc 5.1  --  phu trach: Dung (xem report/PHANCONG.md)
% nguon: rtl/cnn1d/de10_lite_top.v, rtl/cnn1d/top/ecg_accel_top.v, rtl/cnn1d/comm/uart/uart_rx.v, uart_tx.v
\section{Tầng vật lý UART}\label{sec:uartvl}
Host kết nối với board DE10-Lite qua bộ thu phát không đồng bộ vạn năng (Universal Asynchronous Receiver/Transmitter, UART) mức logic $3{,}3$~V thông qua cầu chuyển đổi USB--UART CH340.
Giao diện trên FPGA (\texttt{de10\_lite\_top}) gồm chân thu \texttt{uart\_rx} (\texttt{PIN\_AB5}), chân phát \texttt{uart\_tx} (\texttt{PIN\_AB6}) và đất chung, truyền theo định dạng \emph{8N1} ($1$~bit start mức~$0$, $8$~bit dữ liệu từ bit thấp trước, $1$~bit stop mức~$1$) không dùng điều khiển luồng phần cứng.
Mạch vận hành với xung nhịp hệ thống \texttt{CLK\_FREQ} $=50$~MHz và tốc độ truyền \texttt{BAUD\_RATE} $=1\,500\,000$~baud (hỗ trợ cấu hình ở $921\,600$~baud và $1\,500\,000$~baud), chia chu kỳ bit theo $\texttt{DIV}=\lfloor\texttt{CLK\_FREQ}/\texttt{BAUD\_RATE}\rfloor$ với điều kiện thiết kế $\texttt{CLK\_FREQ}/\texttt{BAUD\_RATE}\ge8$.

Module thu \texttt{uart\_receiver} (\texttt{comm/\allowbreak uart/\allowbreak uart\_rx.v}) đồng bộ tín hiệu \texttt{rx\_in} qua hai tầng thanh ghi liên tiếp (\texttt{rx\_meta}, \texttt{rx\_sync}), sau đó giải mã byte bằng máy trạng thái hữu hạn (Finite State Machine, FSM) năm trạng thái: \texttt{IDLE}, \texttt{START}, \texttt{DATA}, \texttt{STOP} và \texttt{RECOVER}.
Khi gặp cạnh xuống ở \texttt{IDLE}, mạch sang \texttt{START} và đếm nửa chu kỳ bit ($\texttt{DIV}/2-1$) để kiểm tra tại tâm bit start: nếu đường dây về mức~$1$ sớm thì loại bỏ start giả và quay lại \texttt{IDLE}; nếu giữ mức~$0$ thì chuyển sang \texttt{DATA} lấy mẫu $8$~bit dữ liệu cách đều mỗi $\texttt{DIV}$ chu kỳ.
Tại \texttt{STOP}, mạch lấy mẫu tâm bit stop: nếu bằng~$1$ thì xuất \texttt{data\_out} kèm xung \texttt{rx\_ready}; nếu bằng~$0$ thì phát xung \texttt{frame\_error} và đợi ở \texttt{RECOVER} cho tới khi đường dây trở về mức nghỉ~$1$.
Ở chiều ngược lại, \texttt{uart\_transmitter} (\texttt{comm/\allowbreak uart/\allowbreak uart\_tx.v}) nhận byte mới khi cờ \texttt{tx\_ready} bật, dịch nối tiếp bit start, $8$~bit dữ liệu và giữ đủ trọn độ rộng bit stop trước khi hoàn tất khung.
